\usepackage[letterpaper, portrait, margin=1in]{geometry}
\usepackage{algorithm}
\usepackage[noend]{algpseudocode}
\usepackage{url}
\usepackage{amsmath,amssymb,amsthm}
\usepackage[dvipsnames,svgnames,table]{xcolor}
\usepackage{tikz}
\usepackage{tikz-3dplot}
\usepackage{pics/tikz-3dplot-circleofsphere}
\usepackage{listofitems}
\usepackage{thmtools,thm-restate}
\usepackage[noabbrev,capitalise,nameinlink]{cleveref}
\usepackage{mathtools}
\usepackage{xspace}
\usepackage{verbatim}
\usepackage{mathrsfs}
\usepackage{pgf}
%\usepackage[dvipsnames]{xcolor}
\usepackage{todo}
\usepackage{tabularx}
\usepackage{enumitem}
\usepackage{derivative}
\usepackage{bm}
\usepackage{multirow}
\usepackage{diagbox}
\usepackage{nicematrix}
\usepackage[most]{tcolorbox}
\usepackage{caption}



\usetikzlibrary{
    babel,
    arrows,
	calc,
	patterns,
	intersections,
    backgrounds,
    trees,
    mindmap,
    fadings,
	decorations.pathreplacing,
	decorations.pathmorphing,
	decorations.text,
	decorations.markings,
    scopes,
	shapes,
	positioning
}

%-------------------------------------------------------------------------------
% Fonts
%-------------------------------------------------------------------------------
\usepackage{dsfont}

\newcommand*\ie{i.\kern.1em e.,\ }
\newcommand*\eg{e.\kern.1em g.,\ }
\newcommand*\cf{c.\kern.1em f.\ }
\newcommand*\almev{a.\kern.1em e.\ }

% -------------------------------------------
% COLORS
% -------------------------------------------

% iroshizuku color names
\definecolor{ama-iro}{RGB}{0, 158, 243.0}
\definecolor{fuyu-gaki}{RGB}{251, 74, 52}
\definecolor{momiji}{RGB}{245, 70, 111}
\definecolor{hotaru-bi}{RGB}{229,221,58} % yellowish
\definecolor{kon-peki}{RGB}{1,120,217}
\definecolor{shin-kai}{RGB}{77,98,152}
\definecolor{shin-ryoku}{RGB}{1,145,97}
\definecolor{yama-budo}{RGB}{171,14,122}

%\definecolor{theoremcolor}{RGB}{77,98,152}
%\definecolor{examplecolor}{RGB}{210, 229, 205}
%\definecolor{examplecolor}{RGB}{255, 213, 179}
%\definecolor{definitioncolor}{RGB}{255, 213, 179}

\definecolor{citecolor}{HTML}{5E9100}

\definecolor{theoremcolor}{HTML}{FFE1D9}
\definecolor{resultcolor}{HTML}{FFE1D9}
%\definecolor{theoremcolor}{HTML}{C9E1F1}
%\definecolor{resultcolor}{HTML}{C9E1F1}

\definecolor{constraintcolor}{HTML}{BBD49B}
\definecolor{goalcolor}{HTML}{CAE4A7}

\definecolor{remarkcolor}{HTML}{E3EEC7}

\definecolor{definitioncolor}{HTML}{FCDFBE}

\definecolor{examplecolor}{HTML}{F9E9D9}
\definecolor{questioncolor}{HTML}{DDE8EB}
\definecolor{conjecturecolor}{HTML}{DDE8EB}

\definecolor{captioncolor}{RGB}{128,128,128}


\newcommand{\nathan}[1]{{\color{ama-iro} \textbf{NH:} #1}}
\newcommand{\lorenzo}[1]{{\color{shin-ryoku}\textbf{LB:} #1}}


% hyperref colors:


% QED symbol

%--------------------------------------------------------------------------------
% Caption formatting
%--------------------------------------------------------------------------------

\definecolor{captioncolor}{RGB}{128,128,128}

\captionsetup[figure]{labelfont={tt, small, color=momiji}, textfont={tt,small,color=captioncolor},}

%-------------------------------------------------------------------------------
% Theorems
%-------------------------------------------------------------------------------
\topsep=1em % spacing between plain theorems
\theoremstyle{plain}
\newtheorem{theorem}{Theorem}[section]
\newtheorem{lemma}[theorem]{Lemma}
\newtheorem{fact}[theorem]{Fact}
\newtheorem{proposition}[theorem]{Proposition}
\newtheorem{claim}[theorem]{Claim}
\newtheorem{corollary}[theorem]{Corollary}
\newtheorem{question}[theorem]{Open Problem}
\newtheorem{question*}{Question}
\newtheorem{conjecture}[theorem]{Conjecture}
\newtheorem{hypothesis}[theorem]{Hypothesis}
\newtheorem{observation}[theorem]{Observation}
\newtheorem{assumption}[theorem]{Assumption}

\crefname{claim}{Claim}{Claims}
\crefname{fact}{Fact}{Facts}

\theoremstyle{definition}
\newtheorem{condition}[theorem]{Condition}
\newtheorem{definition}[theorem]{Definition}
\newtheorem{remark}[theorem]{Remark}
\newtheorem{example}[theorem]{Example}
\newtheorem{goal}[theorem]{Goal}
\newtheorem{goal*}{Goal}

\theoremstyle{plain}

% --------------------------------------------------------------------------------
% RESTATED THEOREMS
% --------------------------------------------------------------------------------

\newenvironment{rtheorem}[1]{\renewcommand\thetheorem{\ref*{#1}}\begin{theorem}[Restated]}{\end{theorem}}

\newenvironment{rlemma}[1]{\renewcommand\thelemma{\ref*{#1}}\begin{lemma}[Restated]}{\end{lemma}}

\newenvironment{rdefinition}[1]{\renewcommand\thedefinition{\ref*{#1}}\begin{definition}[Restated]}{\end{definition}}

% --------------------------------------------------------------------------------
% BOX THEOREMS
% --------------------------------------------------------------------------------

\newenvironment{boxtheorem}{\begin{theorem}}{\end{theorem}}
\tcolorboxenvironment{boxtheorem}{colback=theoremcolor, colframe=white,
    colbacktitle=theoremcolor, coltitle=theoremcolor}

\newenvironment{rboxtheorem}[1]{\begin{rtheorem}{#1}}{\end{rtheorem}}
\tcolorboxenvironment{rboxtheorem}{colback=theoremcolor, colframe=white,
	colbacktitle=theoremcolor, coltitle=theoremcolor}

\newenvironment{boxlemma}{\begin{lemma}}{\end{lemma}}
\tcolorboxenvironment{boxlemma}{colback=resultcolor, colframe=white,
    colbacktitle=resultcolor, coltitle=resultcolor}

\newenvironment{rboxlemma}[1]{\begin{rlemma}{#1}}{\end{rlemma}}
\tcolorboxenvironment{rboxlemma}{colback=resultcolor, colframe=white,
	colbacktitle=resultcolor, coltitle=resultcolor}

\newenvironment{boxproposition}{\begin{proposition}}{\end{proposition}}
\tcolorboxenvironment{boxproposition}{colback=resultcolor, colframe=white,
    colbacktitle=resultcolor, coltitle=resultcolor}

\newenvironment{boxcorollary}{\begin{corollary}}{\end{corollary}}
\tcolorboxenvironment{boxcorollary}{colback=resultcolor, colframe=white,
    colbacktitle=resultcolor, coltitle=resultcolor}


\newenvironment{boxobservation}{\begin{observation}}{\end{observation}}
\tcolorboxenvironment{boxobservation}{colback=resultcolor, colframe=white,
    colbacktitle=resultcolor, coltitle=resultcolor}

\newenvironment{boxquestion}{\begin{question}}{\end{question}}
\tcolorboxenvironment{boxquestion}{colback=questioncolor, colframe=white,
    colbacktitle=questioncolor, coltitle=shin-kai}
\newenvironment{boxquestion*}{\begin{question*}}{\end{question*}}
\tcolorboxenvironment{boxquestion*}{colback=questioncolor, colframe=white,
    colbacktitle=questioncolor, coltitle=questioncolor}

\newenvironment{boxdefinition}{\begin{definition}}{\end{definition}}
\tcolorboxenvironment{boxdefinition}{colback=definitioncolor, colframe=white,
    colbacktitle=definitioncolor, coltitle=definitioncolor}

\newenvironment{rboxdefinition}[1]{\begin{rdefinition}{#1}}{\end{rdefinition}}
\tcolorboxenvironment{rboxdefinition}{colback=definitioncolor, colframe=white,
	colbacktitle=definitioncolor, coltitle=definitioncolor}

\newenvironment{boxassumption}{\begin{assumption}}{\end{assumption}}
\tcolorboxenvironment{boxassumption}{colback=definitioncolor, colframe=white,
    colbacktitle=definitioncolor, coltitle=definitioncolor}

\newenvironment{boxexample}{\begin{example}}{\end{example}}
\tcolorboxenvironment{boxexample}{colback=examplecolor, colframe=white,
    colbacktitle=examplecolor, coltitle=examplecolor}

\newtheorem{exercise}[theorem]{Exercise}
\newenvironment{boxexercise}{\begin{exercise}}{\end{exercise}}
\tcolorboxenvironment{boxexercise}{colback=exercisecolor, colframe=white,
    colbacktitle=exercisecolor, coltitle=exercisecolor}

\newenvironment{boxgoal}{\begin{goal}}{\end{goal}}
\tcolorboxenvironment{boxgoal}{colback=goalcolor, colframe=white,
    colbacktitle=goalcolor, coltitle=goalcolor}

\newenvironment{boxgoal*}{\begin{goal*}}{\end{goal*}}
\tcolorboxenvironment{boxgoal*}{colback=goalcolor, colframe=white,
    colbacktitle=goalcolor, coltitle=goalcolor}

\newenvironment{boxremark}{\begin{remark}}{\end{remark}}
\tcolorboxenvironment{boxremark}{colback=remarkcolor, colframe=white,
    colbacktitle=remarkcolor, coltitle=remarkcolor}

\newenvironment{boxconjecture}{\begin{conjecture}}{\end{conjecture}}
\tcolorboxenvironment{boxconjecture}{colback=conjecturecolor, colframe=white,
    colbacktitle=conjecturecolor, coltitle=shin-kai}
\newenvironment{boxconjecture*}{\begin{conjecture*}}{\end{conjecture*}}
\tcolorboxenvironment{boxconjecture*}{colback=conjecturecolor, colframe=white,
    colbacktitle=conjecturecolor, coltitle=conjecturecolor}

%-------------------------------------------------------------------------------
% Environments
%-------------------------------------------------------------------------------

\newcommand{\ignore}[1]{}


%-------------------------------------------------------------------------------
% Common functions
%-------------------------------------------------------------------------------
\DeclareMathOperator{\supp}{supp}   % support
\DeclareMathOperator{\poly}{poly}
\DeclareMathOperator{\sign}{sign}
\DeclareMathOperator{\dist}{dist}


\newcommand{\VC}{\mathsf{VC}}

%-------------------------------------------------------------------------------
% Probability stuff
%-------------------------------------------------------------------------------

% Variance:
\DeclareMathOperator{\Cov}{Cov}
%\newcommand{\Var}[1]{\bV \left[ #1 \right]}
\newcommand{\Var}[1]{\mathrm{Var} \left[ #1 \right]}
\newcommand{\Varu}[2]{\underset{ #1 } {\mathrm{Var}} \left[ #2 \right]}
\newcommand{\Varuc}[3]{\underset{ #1 } {\mathrm{Var}} \left[ #2 \;\; \left| \;\; #3 \right. \right]}


% Influence:
\newcommand{\Inf}[1]{\mathrm{Inf} \left[ #1 \right]}


% Expectation:
\newcommand{\Ex}[1]{\bE \left[ #1 \right]}
\newcommand{\Exu}[2]{\underset{#1} \bE \left[ #2 \right] }
\newcommand{\Exuc}[3]{\underset{#1} \bE \left[ #2 \;\; \left| \;\; #3
\right.\right] }

% Probability:
\renewcommand{\Pr}[1]{\bP \left[ #1 \right]} % Probability
\newcommand{\Pru}[2]{\underset{ #1 }\bP \left[ #2 \right]}
\newcommand{\Pruc}[3]{\underset{ #1 }\bP \left[ #2 \;\; \mid \;\; #3 \right]}

%-------------------------------------------------------------------------------
% Common symbols
%-------------------------------------------------------------------------------
% definitions (from https://tex.stackexchange.com/questions/4216/how-to-typeset-correctly)
\newcommand*{\define}{\coloneqq} % Mika: This is the standard symbol that the average reader is used to.

% exponential function
%\renewcommand{\exp}[1]{\mathrm{exp}\left( #1 \right)}

% floor and ceiling
\newcommand{\floor}[1]{\ensuremath{\lfloor #1 \rfloor}}
\newcommand{\ceil}[1]{\ensuremath{\lceil #1 \rceil}}

\newcommand{\inn}[1]{\langle #1 \rangle}

%\newcommand{\abs}[1]{| #1 |}
\DeclarePairedDelimiter{\abs}{\lvert}{\rvert}

\newcommand{\id}{\ensuremath{\mathds{1}}}
\newcommand{\ind}[1]{\mathds{1} \left[ #1 \right] }

% question-mark relations
\newcommand\ltquestion{\stackrel{\mathclap{\normalfont\mbox{\tiny ?}}}{<}}
\newcommand\lequestion{\stackrel{\mathclap{\normalfont\mbox{\tiny ?}}}{\le}}
\newcommand\gtquestion{\stackrel{\mathclap{\normalfont\mbox{\tiny ?}}}{>}}
\newcommand\gequestion{\stackrel{\mathclap{\normalfont\mbox{\tiny ?}}}{\ge}}
\newcommand\eqquestion{\stackrel{\mathclap{\normalfont\mbox{\tiny ?}}}{=}}

\newcommand\lesssimquestion{\stackrel{\mathclap{\normalfont\mbox{?}}}{\lesssim}}
\newcommand\gtrsimquestion{\stackrel{\mathclap{\normalfont\mbox{?}}}{\gtrsim}}
\newcommand\approxquestion{\stackrel{\mathclap{\normalfont\mbox{?}}}{\approx}}

%-------------------------------------------------------------------------------
% Common sets
%-------------------------------------------------------------------------------
\newcommand{\zo}{\{0,1\}}
\newcommand{\pmset}{\{\pm 1\}}

%-------------------------------------------------------------------------------
% Shortcuts for special characters
%-------------------------------------------------------------------------------

% Calligraphic characters
\newcommand{\cA}{\ensuremath{\mathcal{A}}}
\newcommand{\cB}{\ensuremath{\mathcal{B}}}
\newcommand{\cC}{\ensuremath{\mathcal{C}}}
\newcommand{\cD}{\ensuremath{\mathcal{D}}}
\newcommand{\cE}{\ensuremath{\mathcal{E}}}
\newcommand{\cF}{\ensuremath{\mathcal{F}}}
\newcommand{\cG}{\ensuremath{\mathcal{G}}}
\newcommand{\cH}{\ensuremath{\mathcal{H}}}
\newcommand{\cI}{\ensuremath{\mathcal{I}}}
\newcommand{\cJ}{\ensuremath{\mathcal{J}}}
\newcommand{\cK}{\ensuremath{\mathcal{K}}}
\newcommand{\cL}{\ensuremath{\mathcal{L}}}
\newcommand{\cM}{\ensuremath{\mathcal{M}}}
\newcommand{\cN}{\ensuremath{\mathcal{N}}}
\newcommand{\cO}{\ensuremath{\mathcal{O}}}
\newcommand{\cP}{\ensuremath{\mathcal{P}}}
\newcommand{\cQ}{\ensuremath{\mathcal{Q}}}
\newcommand{\cR}{\ensuremath{\mathcal{R}}}
\newcommand{\cS}{\ensuremath{\mathcal{S}}}
\newcommand{\cT}{\ensuremath{\mathcal{T}}}
\newcommand{\cU}{\ensuremath{\mathcal{U}}}
\newcommand{\cV}{\ensuremath{\mathcal{V}}}
\newcommand{\cY}{\ensuremath{\mathcal{Y}}}
\newcommand{\cX}{\ensuremath{\mathcal{X}}}
\newcommand{\cZ}{\ensuremath{\mathcal{Z}}}

% Blackboard Bold characters
\newcommand{\bA}{\ensuremath{\mathbb{A}}}
\newcommand{\bB}{\ensuremath{\mathbb{B}}}
\newcommand{\bC}{\ensuremath{\mathbb{C}}}
\newcommand{\bD}{\ensuremath{\mathbb{D}}}
\newcommand{\bE}{\ensuremath{\mathbb{E}}}
\newcommand{\bF}{\ensuremath{\mathbb{F}}}
\newcommand{\bN}{\ensuremath{\mathbb{N}}}
\newcommand{\bP}{\ensuremath{\mathbb{P}}}
\newcommand{\bR}{\ensuremath{\mathbb{R}}}
\newcommand{\bS}{\ensuremath{\mathbb{S}}}
\newcommand{\bT}{\ensuremath{\mathbb{T}}}
\newcommand{\bV}{\ensuremath{\mathbb{V}}}
\newcommand{\bZ}{\ensuremath{\mathbb{Z}}}

% Bold characters
%\newcommand{\bA}{\boldsymbol{A}}
%\newcommand{\bB}{\boldsymbol{B}}
%\newcommand{\bC}{\boldsymbol{C}}
%\newcommand{\bD}{\boldsymbol{D}}
%\newcommand{\bE}{\boldsymbol{E}}
%\newcommand{\bF}{\boldsymbol{F}}
%\newcommand{\bG}{\boldsymbol{G}}
%\newcommand{\bH}{\boldsymbol{H}}
%\newcommand{\bI}{\boldsymbol{I}}
%\newcommand{\bJ}{\boldsymbol{J}}
%\newcommand{\bK}{\boldsymbol{K}}
%\newcommand{\bL}{\boldsymbol{L}}
%\newcommand{\bM}{\boldsymbol{M}}
%\newcommand{\bN}{\boldsymbol{N}}
%\newcommand{\bO}{\boldsymbol{O}}
%\newcommand{\bP}{\boldsymbol{P}}
%\newcommand{\bQ}{\boldsymbol{Q}}
%\newcommand{\bR}{{\boldsymbol{R}}}
%\newcommand{\bS}{\boldsymbol{S}}
%\newcommand{\bT}{{\boldsymbol{T}}}
%\newcommand{\bU}{\boldsymbol{U}}
%\newcommand{\bV}{\boldsymbol{V}}
%\newcommand{\bW}{\boldsymbol{W}}
%\newcommand{\bX}{\boldsymbol{X}}
%\newcommand{\bY}{\boldsymbol{Y}}
%\newcommand{\bZ}{\boldsymbol{Z}}



\tikzset{
    >=latex,
    graph-vert/.style 2 args = {
        draw,
        color = #1!70!black,
        circle,
        thick,
        inner sep = 0pt,
        minimum size = #2,
        fill = #1!50,
        fill opacity = 0.5
    },
    graph-vert/.default = {blue}{0.15cm},
    triangle/.style 2 args = {
        #1!70!black,
        fill = #1!30,
        fill opacity = #2
    },
    universe-rect/.pic = {
        \fill[black!10, rounded corners = 3pt] (-0.1, 0.1) rectangle (#1 + 0.1, -#1 - 0.1);
        \fill[white] (0, 0) rectangle (#1, -#1);
        \draw[step = 0.5, black] (0, 0) grid (#1, -#1);
        \draw[thick] (-0, 0) rectangle (#1, -#1);
    },
    vert/.style = {
        draw,
        thick,
        rectangle,
        rounded corners = 2pt,
%        ellipse
    },
    semisim/.style = {
        ->,
        blue,
        dashed,
        decorate,
        decoration = {
            snake,
            amplitude = 0.5,
            segment length = 2
        }
    }
}